% exp-fips203-mlkem-pke-encrypt-k3 — FIPS 203 Alg.14 ML-KEM-768 PKE Encrypt
% 编译：bash scripts/xelatex-clean.sh examples/incubating/ml-kem/ml-kem-768/exp-fips203-mlkem-pke-encrypt-k3/exp-fips203-mlkem-pke-encrypt-k3-实现方案-customspec.tex
\documentclass[UTF8,a4paper,11pt]{ctexart}
\usepackage{amsmath,amssymb}
\usepackage{booktabs}
\usepackage{geometry}
\usepackage{hyperref}
\usepackage{longtable}
\usepackage{array}
\usepackage{seqsplit}
\geometry{margin=2.2cm}
\newcolumntype{L}[1]{>{\raggedright\arraybackslash}p{#1}}
\newcommand{\apiname}[1]{\texttt{\seqsplit{#1}}}
\newcommand{\dirname}[1]{\texttt{\seqsplit{#1}}}
\hypersetup{colorlinks=true,linkcolor=blue,urlcolor=blue}

\title{exp-fips203-mlkem-pke-encrypt-k3\\FIPS 203 Alg.14 ML-KEM-768 PKE Encrypt 实现方案}
\author{ascendc 工程（ML-KEM-768 W4a / E14）}
\date{2026-07-26}

\begin{document}
\maketitle

\begin{abstract}
本文档描述 \dirname{examples/incubating/ml-kem/ml-kem-768/exp-fips203-mlkem-pke-encrypt-k3/}：
在 \textbf{Atlas A2 / Ascend910B4} 上以 AscendC 实现 FIPS 203 \textbf{Algorithm 14}，即 ML-KEM-768（$k=3$）\textbf{PKE Encrypt}。

\textbf{参数锁定}：几何、I/O 与 launch 全部锁定到
\dirname{docs/specs/fips203-mlkem768-parameter-card.md} §3.2 D14。
\textbf{行为基线}：活跃探针
\dirname{ascendc-tests/ml-kem/ml-kem-768/pass-fix-f203-alg14-pke-encrypt-device-k3/}。
\textbf{生产 I/O}：\texttt{input/} 为 \texttt{ek\_pke.bin}（1184B）、\texttt{m.bin}（32B）、\texttt{coins.bin}（32B）与静态 LUT；
\texttt{output/} 仅 \texttt{c.bin}（1088B）。
\textbf{计划 AI Core 数}：2 launch；Launch 1 为 AIV\_ONLY \texttt{blockDim=2}，Launch 2 为 MIX \texttt{blockDim=1}（1 AIC + 2 AIV）。
\textbf{禁止}：创建 \dirname{examples/stable/ml-kem/ml-kem-768/}；从任何 \dirname{frozen/} 目录复制、改写或移植源码；用 k4 的 \texttt{ek=1568}、\texttt{c=1568}、\(\hat A[16]\)、\texttt{re[9]} 或 INTT batch 6/8 替代 k3 几何。
\end{abstract}

\tableofcontents
\newpage

\section{工程边界}

\subsection{允许的代码来源与依赖}

\begin{longtable}{@{}L{4.6cm}L{8.6cm}@{}}
\toprule
来源 & 用途 \\
\midrule
\dirname{ascendc-tests/ml-kem/ml-kem-768/pass-fix-f203-alg14-pke-encrypt-device-k3/} & 活跃 D14 k3 行为基线；可复制后自包含适配到本 exp，最终不得保留对探针目录的编译期 include/import \\
\dirname{examples/incubating/ml-kem/ml-kem-1024/exp-fips203-mlkem-pke-encrypt-k4/} & 活跃 k4 incubating 结构参考；仅 retarget 到 k3 锁定几何与字节长 \\
\dirname{library/shared/} & SHAKE/Keccak、公共设备侧头文件、统一整数压缩辅助等共享库 \\
CANN / tikicpulib / CAModel & 编译、CPU 孪生与 SIM 运行环境 \\
\bottomrule
\end{longtable}

\subsection{禁止项}

\begin{itemize}
  \item 禁止从 \dirname{ascendc-tests/frozen/} 或 \dirname{examples/frozen/} 拿源码、customspec 或路线。
  \item 禁止将 \(\hat A\)、\texttt{re}、\(\hat r\)、\(u\)、\(v\)、\(\hat t\) 等中间态作为默认 I/O 落盘；debug dump 只能显式开关启用。
  \item 禁止把 k4 的 \texttt{du=11/dv=5}、\texttt{c1=1408/c2=160}、\texttt{re[9]}、INTT batch 6/8 或零垫 4/8 当作默认。
  \item 禁止新建或写入 stable-768；本目录只作为 incubating 预研实现。
\end{itemize}

\section{数学语义（Alg.14，$k=3$，$n=256$，$q=3329$）}

\subsection{输入输出}

\begin{longtable}{@{}L{3.3cm}L{10.2cm}@{}}
\toprule
张量 / 文件 & 契约 \\
\midrule
\texttt{input/ek\_pke.bin} & 1184B；前 1152B 为 \(\hat t\) 的 \(\mathrm{ByteEncode}_{12}\)，末 32B 为 \(\rho\) \\
\texttt{input/m.bin} & 32B 明文消息；尾段按 \(\mathrm{ByteDecode}_{1}\) 后 \(\mathrm{Decompress}_{1}\) 注入 \(v\) \\
\texttt{input/coins.bin} & 32B Alg.14 随机性；用于 PRF 生成 \(r\Vert e_1\Vert e_2\) \\
\texttt{input/lut\_even\_stacked.bin}、\texttt{lut\_odd\_stacked.bin} & NTT/INTT Stage2 LUT；由本目录 golden 脚本生成；不是用户可见交付输入 \\
\texttt{output/c.bin} & 1088B，\(c=c_1\Vert c_2\)，其中 \(c_1=960\)B（\(d_u=10\)，3 poly），\(c_2=128\)B（\(d_v=4\)，1 poly） \\
\bottomrule
\end{longtable}

\subsection{阶段公式}

\begin{enumerate}
  \item 行 2：从 \texttt{ek\_pke[0:1152]} 解码 \(\hat t\in R_q^3\)，此步骤在 compute launch 内完成，Host 不预解码落盘。
  \item 行 3--7：取 \(\rho=\texttt{ek\_pke[1152:1184]}\)，对 \((j,p)\in\{0,1,2\}^2\) 生成 \(\hat A[j,p]\)，GM 形状锁定为 \texttt{[9,256] int32}。
  \item 行 8--15：用 \(\eta_1=\eta_2=2\) 从 \texttt{coins} 采样 \(r\Vert e_1\Vert e_2\) 共 7 个 poly；GM 形状锁定为 \texttt{[7,256]}。
  \item 行 16--18：计算 \(\hat r=\mathrm{NTT}(r)\)，再得 \(\hat u_j=\sum_p \hat A[j,p]\circ\hat r_p\) 与 \(\widehat{tr}=\sum_p \hat t_p\circ\hat r_p\)。
  \item 行 19--21：对 \(\hat u_0,\hat u_1,\hat u_2,\widehat{tr}\) 执行真 \textbf{INTT batch=4}；得到 \(u_0,u_1,u_2,v\) 后分别加 \(e_1\) 与 \(e_2+\mu(m)\)。
  \item 行 22：\(c_1=\mathrm{ByteEncode}_{10}(\mathrm{Compress}_{10}(u))\)，\(c_2=\mathrm{ByteEncode}_{4}(\mathrm{Compress}_{4}(v))\)，输出 \(c=c_1\Vert c_2\)。
\end{enumerate}

\section{AscendC 实现规划}

\subsection{Launch 与 AI Core 规模}

\begin{longtable}{@{}L{2.5cm}L{3.0cm}L{7.6cm}@{}}
\toprule
Launch & 核类型 / blockDim & 数据划分与理由 \\
\midrule
1 prep & AIV\_ONLY，\texttt{blockDim=2} & \(\hat A\) 共 9 poly 按 5+4 分给两个 AIV；\(r\Vert e_1\Vert e_2\) 共 7 poly 由 block0 生成并按真实 7 poly 写 GM。选择 2 AIV 是参数卡 §3.2 D14 锁定值，继承 D14 探针已验同步与 UB 分配。 \\
2 compute+tail & MIX，\texttt{blockDim=1} & 1 AIC + 2 AIV；NTT(\(r\)) 只处理真实 3 poly；内积输出 4 个待 INTT poly；INTT batch4 的 AIV 映射为连续 2+2。选择单 MIX AI Core 是 D14 锁定生产档，避免跨 AI Core GM 交换。 \\
\bottomrule
\end{longtable}

\subsection{GM 几何}

\begin{longtable}{@{}L{3.7cm}L{9.5cm}@{}}
\toprule
对象 & 锁定几何 \\
\midrule
\(\hat A\) & \texttt{[9,256] int32}；9 个真实 poly，禁止 pad 到 16 \\
\texttt{re} & \texttt{[7,256] int32}；\texttt{r[0..2] || e1[0..2] || e2[0]} \\
\(\hat r\) NTT & polyvec3；复用 B5/NTT 真实 poly 几何，不引入假 poly \\
INTT batch4 & S0 \texttt{[8,256] int8}，Stage2 \texttt{mat\_c [32,128] int32}；Cube 的 M 对齐垫只是硬件对齐，不表示假 poly \\
Tail pack & \(u\) 三 poly 用 \(d_u=10\) pack 为 960B；\(v\) 一 poly 用 \(d_v=4\) pack 为 128B \\
\bottomrule
\end{longtable}

\subsection{同步与流水}

Prep 内部使用 UB/GM 分段搬运与必要的 \apiname{PipeBarrier}；两次 launch 之间由 Host 顺序保证 GM 可见性。
Compute 内部遵守 D14/B5/B6 已验模式：AIV 写入平面 Stage1/Stage3 GM 后以 CrossCore flag 与 AIC 同步；AIC 写完 \texttt{mat\_c} 后 AIV 等待再做 Stage3 merge、内积与 tail pack。
CPU 仅作为逻辑孪生；SIM 是同步与搬运验收门槛。

\section{AscendC API 表}

本实现不引入 E14 独有的新 AscendC API；下表 API 均已有索引记录，使用约束按
\dirname{library/documents/CANN-AscendC算子开发接口参考-查阅索引.md} 执行。

\begin{longtable}{@{}L{3.2cm}L{2.7cm}L{2.5cm}L{2.2cm}L{3.0cm}@{}}
\toprule
API 名 & 分类 / 文档 & Atlas A2 & 本用例 dtype & 备注 \\
\midrule
\apiname{GetBlockIdx} / \apiname{GetSubBlockIdx} & 系统变量 / 资源 & √ & 标量索引 & 区分 AIV 分片与 MIX 子核；禁止把 CPU 打印误认为多核性能 \\
\apiname{GlobalTensor} / \apiname{LocalTensor} & 基础结构 & √ & \texttt{uint8/int8/int32} & GM/UB 张量契约必须与本 spec 几何一致 \\
\apiname{DataCopy} & Memory 数据搬运 & √ & \texttt{uint8/int8/int32} & 连续块搬运；长度与地址遵守 32B 对齐约束 \\
\apiname{PipeBarrier} & Pipe 同步 & √ & — & MTE/Vector/Cube 阶段交界显式同步 \\
\apiname{CrossCoreSetFlag} / \apiname{CrossCoreWaitFlag} & MIX 同步 & √ & — & AIV/AIC 通过 GM 交接时使用；CPU 过不能删 \\
\apiname{Mmad} / \apiname{Fixpipe} & 矩阵计算 & √ & \texttt{int8*int8->int32} & Stage2 MMAD；INTT batch4 的逻辑 M=8，硬件 M 对齐垫到 16 \\
\apiname{ShiftRight} & 矢量算术 & √ & \texttt{int32} & limb 拆分、Barrett 压缩与 ByteEncode 辅助右移 \\
\apiname{Muls} / \apiname{Add} / \apiname{Sub} & 矢量算术 & √ & \texttt{int32/int16} & limb、模加减、压缩/编码辅助 \\
\apiname{Cast} & 类型转换 & √（受限） & \texttt{int32/int16/half/int8} & 遵守索引中 A2 Cast 链约束；不得假设存在 \texttt{int32->int8} 单跳 \\
\apiname{Gather} & 矢量搬运/选择 & √ & \texttt{int32} & 仅允许在 Alg.11 / ByteEncode 等非 NTT S1--S3 段按既有路径使用；NTT 三段内禁用 \\
\apiname{GetValue} / \apiname{SetValue} & LocalTensor 标量访问 & √ & \texttt{uint8/int32} & 仅用于 pack 尾部等标量缺口；不得把 Compares bit 包当逐 lane 布尔 \\
\bottomrule
\end{longtable}

\subsection{评估过但未采用}

\begin{longtable}{@{}L{3.8cm}L{9.2cm}@{}}
\toprule
API / 路线 & 不采用原因 \\
\midrule
\apiname{Compares}/\apiname{GatherMask} compact 链 & Alg.14 E14 不需要新增 reject compact；沿用 D14 已验 pack 路径，不新增比较掩码风险 \\
高阶 \apiname{Matmul<>} & 本路线继承 D14 手写 MMAD + 平面 mat\_c 契约；高阶 Matmul 输出布局与 MIX 交接不作为本轮变量 \\
\apiname{Gather} 用于 NTT S1--S3 & 参数卡与 ML-KEM NTT 定稿禁止；NTT 三段保持平面 mat\_c + 连续 poly 批处理 \\
\bottomrule
\end{longtable}

\section{数学步骤覆盖率}

\begin{longtable}{@{}L{4.2cm}L{4.0cm}L{2.0cm}L{3.0cm}@{}}
\toprule
数学步骤 & 实现方式 / API & 覆盖 & 缺口说明 \\
\midrule
取 \(\rho\)、采样 \(\hat A[9]\) & AIV\_ONLY + XOF/rej + \apiname{DataCopy} & 100\% 设备 & 5+4 分片；无零垫 \\
采样 \(r\Vert e_1\Vert e_2[7]\) & AIV\_ONLY + PRF/CBD & 100\% 设备 & 真实 7 poly；不借 k4 9 poly \\
NTT(\(r\)) 与内积 & MIX AIV/AIC + \apiname{Mmad} & 100\% 设备 & NTT 内无 \apiname{Gather} \\
INTT batch4 & MIX AIV/AIC + Stage3 merge & 100\% 设备 & 真 polyvec4，AIV 连续 2+2 \\
Compress/ByteEncode \(d=10/4\) & 统一整数压缩 + pack & 部分向量 & pack 尾部标量缺口沿用 D14/k4 已验实现 \\
\bottomrule
\end{longtable}

\section{验证计划}

\subsection{默认验收}

\begin{verbatim}
cd examples/incubating/ml-kem/ml-kem-768/exp-fips203-mlkem-pke-encrypt-k3
bash run.sh -r cpu -v Ascend910B4
SIM_DIRECT=1 bash run.sh -r sim -v Ascend910B4
\end{verbatim}

期望：\texttt{c.bin} 为 1088B；与 golden 对拍 \texttt{[cmp] c max=0}。
SIM 通过后记录 \texttt{Total tick} 到 \texttt{qa/active\_sim\_regress\_summary.md}，并检查用例根目录无 stray \texttt{core*.dump}、\texttt{profile\_*log*.toml}。

\subsection{可选对照}

后续可补 PKE round-trip 或 liboqs encrypt/decrypt 夹具，但 E14 本轮门禁为 CPU + \texttt{SIM\_DIRECT=1} sim。
liboqs 仅作为黑盒 I/O oracle；禁止把 liboqs 源码实现同构移植进 AscendC。

\section{产出物}

\begin{itemize}
  \item \dirname{exp-fips203-mlkem-pke-encrypt-k3-实现方案-customspec.tex}
  \item \dirname{exp-fips203-mlkem-pke-encrypt-k3-实现方案-customspec.pdf}
  \item 后续实现文件须与本 customspec 同目录自包含，且自研 Python/C/AscendC 写详细中文注释。
\end{itemize}

\end{document}
